\chapter{Inapprossimabilità}

In questo capitolo presentiamo il legame tra Teorema PCP e
(in)approssimabilità di problemi di ottimizzazione combinatoria.
Cominciamo definendo cos'è un algoritmo di approssimazione
e cosa s'intende per inapprossimabilità.
Introduciamo poi i \emph{Constraint Satisfaction Problem} ($\CSP$) e le
relative notazioni che torneranno utili nel Capitolo \ref{chap:pcpproof},
dopodiché diamo una formulazione equivalente del Teorema PCP instrinsecamente
legata all'inapprossimabilità di $\CSP$.
Infine, a partire dal Teorema PCP, dimostriamo risultati di
inapprossimabilità per \texttt{MAX-CLIQUE} e \texttt{MAX-3SAT}.

\section{Approssimare problemi di ottimizzazione}

Un problema di ottimizzazione combinatoria chiede di trovare la soluzione
migliore possibile tra tutte le soluzioni ammissibili, ovvero quella che
massimizza (o minimizza) una determinata funzione obiettivo.
Ad ogni problema di ottimizzazione è naturally associata una versione
decisionale del tipo: ``esiste una soluzione di valore almeno $k$?''.
Se il problema di decisione associato a un problema di ottimizzazione
è $\NP$-completo\footnote{O in generale, non si può risolvere
efficientemente a meno che $\P = \NP$.}, allora calcolare la soluzione ottima
in tempo polinomiale è impossibile, a meno che $\P = \NP$.
È naturale chiedersi se esistono algoritmi efficienti in grado di
trovare una soluzione abbastanza buona, entro un certo fattore dell'ottimo.

\begin{definition}
	Fissato un problema di ottimizzazione, un algoritmo che restituisce sempre
	una soluzione entro un fattore $\alpha$ dall'ottimo si dice
	\emph{$\alpha$-approssimazione}.
	Il fattore di approssimazione $\alpha$ può essere funzione dell'input.
\end{definition}

Consideriamo per esempio \texttt{MAX-SAT}, cioè la versione di ottimizzazione
di \texttt{SAT} in cui l'obiettivo è massimizzare il numero di clausole
soddisfatte.
Consideriamo un assegnamento casuale delle variabili: ogni clausola è
soddisfatta con probabilità almeno $\frac12$, quindi per linearità del
valore atteso ci aspettiamo che almeno metà delle clausole siano soddisfatte.
Sorvoliamo un momento sul fatto che l'algoritmo ``scegli un assegnamento
casuale'' è probabilistico\footnote{È infatti molto facile derandomizzarlo,
cioè ottenerne una controparte deterministica.}.
Un assegnamento che soddisfa sempre metà delle clausole soddisfa a maggior
ragione metà delle clausole soddisfatte da una soluzione ottima, quindi
\texttt{MAX-SAT} si può $\frac12$-approssimare efficientemente.

% TBD. Decidere se citare 7/8 di Hastad

Si parla di \emph{inapprossimabilità} quando, per un dato problema di
ottimizzazione, esiste un limite intrinseco al fattore di approssimazione che
un algoritmo efficiente può garantire, tipicamente a meno che $\P = \NP$.

È bene osservare che la sola $\NP$-completezza non esclude la possibilità di
approssimare molto, magari arbitrariamente, bene la soluzione ottima.
Quasi tutte le dimostrazioni di $\NP$-completezza infatti passano per il
Teorema di Cook-Levin e allo stesso tempo è molto facile trovare assegnamenti
che soddisfano tutte le clausole meno una per le istanze di \texttt{SAT}
prodotte dalla dimostrazione di Cook-Levin\footnote{Il Teorema di Cook-Levin
si dimostra codificando la domanda ``esiste una computazione corrispondente alla
verifica di un certificato accettato?'' in un'istanza di \texttt{SAT}.
La parte ``accettato'' è espressa da una sola clausola, dunque una soluzione
che soddisfa tutte le clausole meno una è data dalla computazione
corrispondente alla verifica di un qualunque certificato, anche non valido.}.
Ciò suggerisce che per dimostrare forti risultati di inapprossimabilità,
sono necessari strumenti diversi. In questo senso il Teorema PCP è stato
fondamentale \cite{Tre04, FGLSS96}.

\section{Teorema PCP e Constraint Satisfaction Problems}

Introduciamo i \emph{Constraint Satisfaction Problem}, cioè una classe di
problemi che generalizza naturalmente \texttt{3SAT} e i problemi di colorazione di
grafi. Fissata un'arietà $k$ e un alfabeto $\Sigma$ definiamo $k\CSP_\Sigma$ il
seguente problema computazionale. Data una collezione $((C_j, R_j))_{j \in [m]}$
di $k$-uple $C_j$ di variabili associate a relazioni $R_j \subseteq \Sigma^k$,
determinare se esiste in $\Sigma^n$ un assegnamento delle variabili che soddisfa
\emph{tutti} i vincoli, cioè tale che per ogni $k$-upla $C_j$ di variabili,
il corrispondente assegnamento appartiene a $R_j$.
Indicheremo con $k\CSP_W$ il problema $k\CSP_{[W]}$.
Secondo la notazione appena introdotta \texttt{3COL}
corrisponde alla restrizione di $2\CSP_3$ alle istanze i cui vincoli sono
espressi da $R_{\text{3col}} = \{(1, 2), (1, 3), (2, 3)\}$.
Ogni variabile corrisponde a un nodo e assume un valore in $\{1, 2, 3\}$; per
ogni coppia di nodi adiacenti si impone la relazione binaria ``$a \neq b$'' tra
le corrispondenti variabili.

Data un'istanza $\varphi$ di $\CSP$, definiamo il suo \emph{valore}
$\val(\varphi)$ come la massima frazione di vincoli soddisfacibili da un
assegnamento delle variabili di $\varphi$. Dunque $\val(\varphi) = 1$ se
$\varphi$ è soddisfacibile e $\val(\varphi) < 1$ altrimenti. Definiamo invece
il \emph{gap} di $\varphi$ come $1 - \val(\varphi)$.

Siamo ora pronti per enunciare la versione di inapprossimabilità del
Teorema PCP.

\begin{theorem}[Teorema PCP -- versione di inapprossimabilità] \label{thm:pcpapprox}
	Esistono costanti $q \in \N$ e $\rho \in (0, 1)$ tali che, per ogni
	$L \in \NP$, esiste una funzione $f$ calcolabile efficientemente che mappa
	istanze di $L$ in istanze di $q\CSP_2$ che soddisfano le seguenti proprietà:
	\begin{itemize}
		\item \textbf{completezza}: se $x \in L$, allora $\val(f(x)) = 1$;
		\item \textbf{robustezza}: se $x \notin L$, allora $\val(f(x)) \le \rho$.
	\end{itemize}
\end{theorem}

Come immediata conseguenza, per ogni $\varepsilon > \rho$ non esiste una
$\varepsilon$-approssimazione efficiente di $\texttt{MAX-}q\CSP_2$
a meno che $\P = \NP$.
Un tale algoritmo, infatti, potrebbe essere utilizzato per decidere
efficientemente l'appartenenza a qualunque linguaggio $L \in \NP$ nel
seguente modo.
Data un'istanza $x$, si esegue l'algoritmo di approssimazione su $f(x)$.
Se $x \in L$, allora $\val(f(x)) = 1$, dunque l'approssimazione restituisce una
soluzione di valore almeno $\varepsilon \val(f(x)) = \varepsilon > \rho$.
Se invece $x \notin L$, allora ogni assegnamento ha valore al più $\rho$.
Di conseguenza, sulla base del valore della soluzione restituita dall'ipotetico
algoritmo di approssimazione, si potrebbe decidere efficientemente
l'appartenenza ad $L$, da cui $\P = \NP$.

Dimostriamo quindi l'equivalenza tra il Teorema \ref{thm:pcp} e il
Teorema \ref{thm:pcpapprox}. Nel Capitolo \ref{chap:pcpproof} useremo
quest'ultima formulazione per dimostrare il Teorema PCP.

\begin{proof}[Dimostrazione di equivalenza]
	$(\implies)$
	Come già fatto durante la dimostrazione del Teorema \ref{thm:exppcp},
	è sufficiente mostrare che la conclusione del Teorema \ref{thm:pcpapprox}
	vale per un fissato linguaggio $\NP$-completo $L$.

	Supponiamo dunque che valga il Teorema \ref{thm:pcp}, cioè che
	$\NP = \PCP(\log n, 1)$. Allora esistono costanti $c, q > 0$ e un
	verificatore probabilistico $V$ che, data un'istanza $x \in \{0, 1\}^n$ di
	$L$, utilizza al più $c \log n$ bit casuali e legge al più $q$ bit
	di una dimostrazione $\pi$.
	Costruiamo l'istanza $\varphi = f(x)$ di $q\CSP_2$ nel modo seguente.
	Le variabili di $\varphi$ corrispondono alle posizioni della dimostrazione
	$\pi$.
	Per ogni stringa casuale $r \in \{0, 1\}^{c \log n}$, si pone un vincolo
	$V_{x,r}$ soddisfatto se e soltanto se $V$ accetta $\pi$ con input $x$ e
	stringa casuale $r$. Poiché $V_{x, r}$ dipende da al più $q$ variabili,
	restringiamo $\varphi$ alle sole variabili che appaiono in uno dei vincoli
	per ottenere un'istanza di taglia polinomiale e computabile con un numero
	polinomiale di simulazioni di verifica.

	Assegnamenti delle variabili di $\varphi$ corrispondono a certificati
	$\pi$, quindi il valore di $\varphi$ corrisponde esattamente alla massima
	probabilità che un certificato sia accettato.
	Le proprietà di completezza e robustezza del Teorema \ref{thm:pcp} si
	traducono allora nelle proprietà di robustezza e completezza per il Teorema
	\ref{thm:pcpapprox} con $\rho = \frac12$.

	$(\impliedby)$
	Viceversa, mostriamo come costruire un valido verificatore $\PCP$ a partire
	dalla trasformazione del Teorema \ref{thm:pcpapprox}.
	Dato input $x$ per un problema $L$, il verificatore lo trasforma nell'istanza
	$f(x)$ di $q\CSP_2$. Dimostrazioni di ``$x \in L$'' corrispondono a validi
	assegnamenti per $f(x)$. Per verificare una dimostrazione, il verificatore
	sceglie un vincolo uniformemente a caso impiegando $O(\log n)$ bit random,
	dopodiché legge le $q$ variabili da cui dipende il vincolo e accetta se e
	soltanto se questo è soddisfatto.
	In virtù del Teorema \ref{thm:pcpapprox}, questo sistema di verifica
	soddisfa la proprietà di completezza e ha robustezza $\rho$.
	È sufficiente ripetere la verifica $O(1)$ volte per ottenere robustezza
	$\frac12$ e dimostrare quindi il Teorema \ref{thm:pcp}.
\end{proof}

\section{Risultati di inapprossimabilità}

\subsection{Inapprossimabilità di \texttt{MAX-CLIQUE}}

%%% BEGIN HERE

% define max-clique

\begin{theorem}[\cite{FGLSS96}] \label{thm:max-clique}
	Per ogni $\varepsilon > 0$, non esiste una $\varepsilon$-approssimazione
	di \texttt{MAX-CLIQUE}, a meno che $\P = \NP$.
\end{theorem}

\begin{proof}
\iffalse
{\color{magenta}%
Dimostriamo il teorema tramite la costruzione del grafo di Feige, Goldwasser, Lovász, Safra e Szegedy \cite{FGLSS96} (grafo FGLSS), riducendo il problema di decidere la soddisfacibilità di un'istanza di $q\CSP_2$ a \texttt{MAX-CLIQUE}.

Sia $L \in \NP$. Per il Teorema \ref{thm:pcpapprox}, esiste un'istanza $\varphi$ di $q\CSP_2$ formata da $m$ vincoli $C_1, \dots, C_m$ sulle variabili $x_1, \dots, x_n \in \{0, 1\}$ tale che $\val(\varphi) = 1$ se $x \in L$ e $\val(\varphi) \le \frac12$ se $x \notin L$.

Costruiamo un grafo non orientato $G = (V, E)$ nel seguente modo:
\begin{itemize}
	\item \textbf{Vertici}: l'insieme $V$ è formato da tutte le coppie $(j, \alpha)$ dove $j \in [m]$ indica l'indice di un vincolo $C_j$ di $\varphi$, e $\alpha \in \{0, 1\}^q$ è un assegnamento locale alle $q$ variabili da cui dipende $C_j$ che soddisfa il vincolo $C_j$. Notiamo che $|V| \le m 2^q = O(m)$.
	\item \textbf{Archi}: due vertici distinti $(j, \alpha)$ e $(j', \alpha')$ sono collegati da un arco in $E$ se e solo se $j \neq j'$ e gli assegnamenti locali $\alpha$ e $\alpha'$ sono \emph{compatibili}, ossia assegnano lo stesso valore a ciascuna variabile condivisa dai vincoli $C_j$ e $C_{j'}$.
\end{itemize}

Mostriamo che la dimensione della massima cricca di $G$, indicata con $\omega(G)$, soddisfa la relazione
\[
	\omega(G) = m \cdot \val(\varphi).
\]
\begin{itemize}
	\item Se $\val(\varphi) \ge \gamma$, allora esiste un assegnamento globale $\sigma \in \{0, 1\}^n$ alle variabili di $\varphi$ che soddisfa almeno $\gamma m$ vincoli, diciamo $C_{j_1}, \dots, C_{j_k}$ con $k \ge \gamma m$. Per ogni vincolo soddisfatto $C_{j_i}$, la restrizione di $\sigma$ alle sue variabili definisce un assegnamento locale soddisfacente $\alpha_i$. I vertici $(j_1, \alpha_1), \dots, (j_k, \alpha_k)$ provengono tutti dal medesimo assegnamento globale $\sigma$, quindi sono a due a due compatibili e appartengono a vincoli distinti. Di conseguenza, costituiscono una cricca di dimensione $k \ge \gamma m$ in $G$, da cui $\omega(G) \ge m \cdot \val(\varphi)$.
	\item Viceversa, sia $S = \{(j_1, \alpha_1), \dots, (j_k, \alpha_k)\}$ una cricca in $G$ di dimensione $k$. Poiché non vi sono archi tra vertici associati allo stesso vincolo, gli indici $j_1, \dots, j_k$ sono tutti distinti. Inoltre, poiché i vertici di $S$ sono mutualmente compatibili, gli assegnamenti locali $\alpha_1, \dots, \alpha_k$ non discordano su alcuna variabile comune e possono essere estesi a un assegnamento globale $\sigma \in \{0, 1\}^n$ che soddisfa contemporaneamente i $k$ vincoli $C_{j_1}, \dots, C_{j_k}$. Ne deriva che $k \le m \cdot \val(\varphi)$, e dunque $\omega(G) \le m \cdot \val(\varphi)$.
\end{itemize}

Dalla relazione $\omega(G) = m \cdot \val(\varphi)$ segue che:
\begin{itemize}
	\item se $x \in L$, allora $\val(\varphi) = 1$ e di conseguenza $\omega(G) = m$;
	\item se $x \notin L$, allora $\val(\varphi) \le \frac12$ e di conseguenza $\omega(G) \le \frac{m}{2}$.
\end{itemize}

Per estendere il risultato da un fattore di inapprossimabilità $\frac12$ a un generico $\varepsilon > 0$, consideriamo il prodotto di grafi (ovvero l'esecuzione di $k$ verifiche indipendenti). Per ogni intero costante $k \ge 1$, costruiamo il grafo $G^k$, i cui vertici sono le $k$-uple di vertici di $G$ e dove due $k$-uple sono adiacenti se e solo se le loro componenti sono uguali o adiacenti in $G$. La dimensione della massima cricca in $G^k$ è $\omega(G^k) = (\omega(G))^k$.
Pertanto, se $x \in L$ si ha $\omega(G^k) = m^k$, mentre se $x \notin L$ si ha $\omega(G^k) \le (m/2)^k = m^k / 2^k$.

Scegliendo l'intero costante $k = \lceil \log_2(1/\varepsilon) \rceil$, la frazione tra il valore nel caso $x \notin L$ e nel caso $x \in L$ diventa $\le 2^{-k} \le \varepsilon$. Poiché $k$ è costante, il grafo $G^k$ ha dimensione $|V|^k = \mathrm{poly}(n)$ e può essere costruito in tempo polinomiale. Se esistesse un algoritmo di $\varepsilon$-approssimazione per \texttt{MAX-CLIQUE}, esso permetterebbe di distinguere il caso $\omega(G^k) = m^k$ dal caso $\omega(G^k) \le \varepsilon m^k$, decidendo l'appartenenza di $x$ a $L$ in tempo polinomiale, il che implicherebbe $\P = \NP$.
}
\fi
\end{proof}

%%% END HERE

\subsection{Inapprossimabilità di \texttt{MAX-3SAT}}

% citare PY91 e spiegare un po' la situazione.

\begin{theorem}
	Esiste un $\varepsilon_1 > 0$ tale che non esiste una
	$\varepsilon_1$-approssimazione di \texttt{MAX-3SAT}, a meno che
	$\P = \NP$.
\end{theorem}

\begin{proof}
\iffalse
	{\color{magenta}%
	Mostriamo una riduzione polinomiale da $q\CSP_2$ a \texttt{MAX-3SAT} che preserva il gap di soddisfacibilità.
	Sia $\varphi$ un'istanza di $q\CSP_2$ formata da $m$ vincoli $C_1, \dots, C_m$ sulle variabili $x_1, \dots, x_n \in \{0, 1\}$, derivante dal Teorema PCP \ref{thm:pcpapprox}.

	Ogni vincolo $C_j$ dipende da al più $q$ variabili originali. Possiamo esprimere la relazione ammissibile per $C_j$ in forma normale congiuntiva (CNF) con al più $2^q$ clausole su tali $q$ variabili. Successivamente, ogni clausola contenente $k$ letterali (con $k \le q$) viene convertita in un insieme di clausole con esattamente 3 letterali mediante l'introduzione di al più $q$ variabili ausiliarie.
	In questo modo, ciascun vincolo $C_j$ viene tradotto in un insieme $S_j$ di clausole 3-CNF definite sulle $q$ variabili originali di $C_j$ e su un insieme di variabili ausiliarie $Y_j$. Indichiamo con $c_q$ il numero totale di clausole 3-CNF appartenenti a $S_j$. Poiché $q = O(1)$ è una costante, anche $c_q = O(1)$ è una costante.

	La riduzione garantisce le seguenti proprietà strutturali per ogni vincolo $C_j$:
	\begin{enumerate}[label=(\roman*)]
		\item Se un assegnamento alle variabili originali soddisfa $C_j$, allora esiste un'estensione alle variabili ausiliarie $Y_j$ che soddisfa \emph{tutte} le $c_q$ clausole di $S_j$.
		\item Se un assegnamento alle variabili originali viola $C_j$, allora per \emph{qualunque} scelta delle variabili ausiliarie $Y_j$, al più $c_q - 1$ clausole di $S_j$ risultano soddisfatte (cioè almeno una clausola di $S_j$ viene violata).
	\end{enumerate}

	Sia $\phi'$ l'istanza di \texttt{MAX-3SAT} formata dalla congiunzione di tutti gli insiemi di clausole $S_1, \dots, S_m$. Il numero totale di clausole di $\phi'$ è $M = m \cdot c_q$.

	Valutiamo il valore $\val(\phi')$ nei due casi:
	\begin{itemize}
		\item \textbf{Completezza}: se $\val(\varphi) = 1$, esiste un assegnamento globale alle variabili $x_1, \dots, x_n$ che soddisfa tutti gli $m$ vincoli $C_j$. Estendendo tale assegnamento in modo opportuno sulle variabili ausiliarie $Y_j$ per ciascun $j \in [m]$, tutte le $c_q$ clausole di ogni $S_j$ risultano soddisfatte. Di conseguenza, tutte le $M$ clausole di $\phi'$ sono soddisfatte e $\val(\phi') = 1$.
		\item \textbf{Robustezza}: se $\val(\varphi) \le \frac12$, allora qualsiasi assegnamento alle variabili $x_1, \dots, x_n$ viola almeno $\frac{m}{2}$ vincoli di $\varphi$. Per ciascun vincolo $C_j$ violato, qualunque assegnamento alle variabili ausiliarie $Y_j$ viola almeno una clausola in $S_j$. Ne consegue che in $\phi'$ risultano violate almeno $\frac{m}{2}$ clausole su $M$, pertanto:
		\[
			\val(\phi') \le \frac{M - m/2}{M} = 1 - \frac{m/2}{m c_q} = 1 - \frac{1}{2 c_q}.
		\]
	\end{itemize}

	Ponendo $\varepsilon_1 = \frac{1}{2 c_q} > 0$, un algoritmo di $\varepsilon_1$-approssimazione per \texttt{MAX-3SAT} (ovvero in grado di garantire soluzioni di valore entro un fattore $(1 - \varepsilon_1)$ dall'ottimo) permetterebbe di distinguere un'istanza con $\val(\phi') = 1$ da un'istanza con $\val(\phi') \le 1 - \varepsilon_1$, decidendo l'appartenenza a $L$ in tempo polinomiale. Si conclude che \texttt{MAX-3SAT} non ammette una $\varepsilon_1$-approssimazione a meno che $\P = \NP$.
}
\fi
\end{proof}
